\subsection{Exact finite remainder bound and affine raw-jet transport}\label{endpoint-relation_balance_original--exact-finite-remainder-bound-and-affine-raw-jet-transport}

This note proves the finite coefficient estimate requested for the volume-comparator calculation. All roots are counted with their original multiplicities. The degree, affine unit, complex phase, domains, and codomains remain explicit throughout.

\subsubsection{1. Polynomial spaces and the remainder map}\label{endpoint-relation_balance_original--polynomial-spaces-and-the-remainder-map}

Let ENDPOINTSOURCE12B6490A442BF168E00ED3A2SLOT0END be an integer and write

ENDPOINTSOURCE12B6490A442BF168E00ED3A2SLOT1END

Repeated entries of ENDPOINTSOURCE12B6490A442BF168E00ED3A2SLOT2END record root multiplicity. Set

ENDPOINTSOURCE12B6490A442BF168E00ED3A2SLOT3END

The exact linear map in this note is

ENDPOINTSOURCE12B6490A442BF168E00ED3A2SLOT4END

Existence follows by subtracting the leading coefficient times the corresponding monomial multiple of the monic polynomial ENDPOINTSOURCE12B6490A442BF168E00ED3A2SLOT5END, decreasing the degree at each step until it is below ENDPOINTSOURCE12B6490A442BF168E00ED3A2SLOT6END. For uniqueness, the difference of two decompositions satisfies ENDPOINTSOURCE12B6490A442BF168E00ED3A2SLOT7END, where the right side has degree below ENDPOINTSOURCE12B6490A442BF168E00ED3A2SLOT8END. A nonzero left side has degree at least ENDPOINTSOURCE12B6490A442BF168E00ED3A2SLOT9END, so both sides vanish. Applying this uniqueness to linear combinations proves linearity. In particular ENDPOINTSOURCE12B6490A442BF168E00ED3A2SLOT10END for ENDPOINTSOURCE12B6490A442BF168E00ED3A2SLOT11END.

\subsubsection{2. The complete homogeneous division formula}\label{endpoint-relation_balance_original--the-complete-homogeneous-division-formula}

Let ENDPOINTSOURCE12B6490A442BF168E00ED3A2SLOT12END be the elementary symmetric polynomial of degree ENDPOINTSOURCE12B6490A442BF168E00ED3A2SLOT13END in the listed roots, with ENDPOINTSOURCE12B6490A442BF168E00ED3A2SLOT14END, and define the complete homogeneous polynomials by

ENDPOINTSOURCE12B6490A442BF168E00ED3A2SLOT15END

As an identity of formal power series,

ENDPOINTSOURCE12B6490A442BF168E00ED3A2SLOT16END

There is no convergence requirement: every coefficient uses only finitely many factors and terms. Thus, for every integer ENDPOINTSOURCE12B6490A442BF168E00ED3A2SLOT17END,

ENDPOINTSOURCE12B6490A442BF168E00ED3A2SLOT18END

For ENDPOINTSOURCE12B6490A442BF168E00ED3A2SLOT19END, put

ENDPOINTSOURCE12B6490A442BF168E00ED3A2SLOT20END

The coefficient of ENDPOINTSOURCE12B6490A442BF168E00ED3A2SLOT21END in ENDPOINTSOURCE12B6490A442BF168E00ED3A2SLOT22END is one. For ENDPOINTSOURCE12B6490A442BF168E00ED3A2SLOT23END, the coefficient of ENDPOINTSOURCE12B6490A442BF168E00ED3A2SLOT24END is the left side of (RB.6), hence zero. Therefore

ENDPOINTSOURCE12B6490A442BF168E00ED3A2SLOT25END

The uniqueness in (RB.3) now gives the exact division identity

ENDPOINTSOURCE12B6490A442BF168E00ED3A2SLOT26END

For completeness its individual coefficients are

ENDPOINTSOURCE12B6490A442BF168E00ED3A2SLOT27END

This expression remains valid for repeated and zero roots without taking limits.

\subsubsection{3. Explicit coefficient-norm bound}\label{endpoint-relation_balance_original--explicit-coefficient-norm-bound}

The coefficient norm is submultiplicative. Indeed, if ENDPOINTSOURCE12B6490A442BF168E00ED3A2SLOT28END and ENDPOINTSOURCE12B6490A442BF168E00ED3A2SLOT29END, then

ENDPOINTSOURCE12B6490A442BF168E00ED3A2SLOT30END

Consequently

ENDPOINTSOURCE12B6490A442BF168E00ED3A2SLOT31END

The number of weak compositions of ENDPOINTSOURCE12B6490A442BF168E00ED3A2SLOT32END into ENDPOINTSOURCE12B6490A442BF168E00ED3A2SLOT33END nonnegative parts is ENDPOINTSOURCE12B6490A442BF168E00ED3A2SLOT34END: place ENDPOINTSOURCE12B6490A442BF168E00ED3A2SLOT35END separators among a row of ENDPOINTSOURCE12B6490A442BF168E00ED3A2SLOT36END markers and those separators. Thus (RB.4) yields

ENDPOINTSOURCE12B6490A442BF168E00ED3A2SLOT37END

In (RB.8), the leading monomial ENDPOINTSOURCE12B6490A442BF168E00ED3A2SLOT38END and the remainder have disjoint degree supports. In particular

ENDPOINTSOURCE12B6490A442BF168E00ED3A2SLOT39END

Combining (RB.10)--(RB.13) proves, with the original radius retained,

ENDPOINTSOURCE12B6490A442BF168E00ED3A2SLOT40END

When ENDPOINTSOURCE12B6490A442BF168E00ED3A2SLOT41END, the finite sum of binomial coefficients is

ENDPOINTSOURCE12B6490A442BF168E00ED3A2SLOT42END

To verify this identity, its ENDPOINTSOURCE12B6490A442BF168E00ED3A2SLOT43END case is ENDPOINTSOURCE12B6490A442BF168E00ED3A2SLOT44END, and its induction step is Pascal's identity ENDPOINTSOURCE12B6490A442BF168E00ED3A2SLOT45END. The right side of (RB.15) increases with ENDPOINTSOURCE12B6490A442BF168E00ED3A2SLOT46END, as is also immediate from its sum expression. Define the exact integer

ENDPOINTSOURCE12B6490A442BF168E00ED3A2SLOT47END

Then ENDPOINTSOURCE12B6490A442BF168E00ED3A2SLOT48END, with ENDPOINTSOURCE12B6490A442BF168E00ED3A2SLOT49END, and every column of the remainder map satisfies

ENDPOINTSOURCE12B6490A442BF168E00ED3A2SLOT50END

The two central coefficients ENDPOINTSOURCE12B6490A442BF168E00ED3A2SLOT51END and ENDPOINTSOURCE12B6490A442BF168E00ED3A2SLOT52END are equal. Their sum is at most the sum of all coefficients of ENDPOINTSOURCE12B6490A442BF168E00ED3A2SLOT53END at ENDPOINTSOURCE12B6490A442BF168E00ED3A2SLOT54END, which is ENDPOINTSOURCE12B6490A442BF168E00ED3A2SLOT55END. Therefore

ENDPOINTSOURCE12B6490A442BF168E00ED3A2SLOT56END

For an arbitrary ENDPOINTSOURCE12B6490A442BF168E00ED3A2SLOT57END, linearity and the triangle inequality give

ENDPOINTSOURCE12B6490A442BF168E00ED3A2SLOT58END

The non-strict version with ENDPOINTSOURCE12B6490A442BF168E00ED3A2SLOT59END holds also for ENDPOINTSOURCE12B6490A442BF168E00ED3A2SLOT60END. In fact the induced operator norm between these coefficient ENDPOINTSOURCE12B6490A442BF168E00ED3A2SLOT61END spaces equals the maximum column norm: the displayed triangle inequality proves the upper bound, and applying the operator to a monomial attaining the finite maximum proves the reverse inequality.

\subsubsection{4. Raw jets and the exact quotient section}\label{endpoint-relation_balance_original--raw-jets-and-the-exact-quotient-section}

Write the distinct roots as ENDPOINTSOURCE12B6490A442BF168E00ED3A2SLOT62END, with multiplicities ENDPOINTSOURCE12B6490A442BF168E00ED3A2SLOT63END, so ENDPOINTSOURCE12B6490A442BF168E00ED3A2SLOT64END. Define the raw-jet space and map by

ENDPOINTSOURCE12B6490A442BF168E00ED3A2SLOT65END

The exact Taylor identity at a root is

ENDPOINTSOURCE12B6490A442BF168E00ED3A2SLOT66END

It shows that the first ENDPOINTSOURCE12B6490A442BF168E00ED3A2SLOT67END raw derivatives vanish precisely when ENDPOINTSOURCE12B6490A442BF168E00ED3A2SLOT68END divides ENDPOINTSOURCE12B6490A442BF168E00ED3A2SLOT69END. Simultaneous divisibility by these distinct-root factors is equivalent to divisibility by their product ENDPOINTSOURCE12B6490A442BF168E00ED3A2SLOT70END: in the factorization of a nonzero polynomial over ENDPOINTSOURCE12B6490A442BF168E00ED3A2SLOT71END, the exponent at each distinct root must be at least the listed multiplicity. The zero polynomial is divisible by every factor. Hence, for ENDPOINTSOURCE12B6490A442BF168E00ED3A2SLOT72END,

ENDPOINTSOURCE12B6490A442BF168E00ED3A2SLOT73END

The restriction ENDPOINTSOURCE12B6490A442BF168E00ED3A2SLOT74END has zero kernel by degree, and both its domain and codomain have dimension ENDPOINTSOURCE12B6490A442BF168E00ED3A2SLOT75END. It is therefore an isomorphism. Equation (RB.3) gives

ENDPOINTSOURCE12B6490A442BF168E00ED3A2SLOT76END

Thus this remainder map preserves every specified raw jet, including every derivative order required by multiplicity. Equivalently, the inclusion of ENDPOINTSOURCE12B6490A442BF168E00ED3A2SLOT77END followed by the quotient map is an isomorphism onto ENDPOINTSOURCE12B6490A442BF168E00ED3A2SLOT78END, and (RB.3) is the unique representative of degree below ENDPOINTSOURCE12B6490A442BF168E00ED3A2SLOT79END in that quotient class.

\subsubsection{\texorpdfstring{5. Original ENDPOINTSOURCE12B6490A442BF168E00ED3A2SLOT80END variable, affine unit, and phase}{5. Original  variable, affine unit, and phase}}\label{endpoint-relation_balance_original--original-s-variable-affine-unit-and-phase}

Let the original monic annihilator be

ENDPOINTSOURCE12B6490A442BF168E00ED3A2SLOT81END

Use the requested transformed-polynomial notation with its full definition:

ENDPOINTSOURCE12B6490A442BF168E00ED3A2SLOT82END

Here the bar symbol is defined by (RB.24); complex conjugation of a polynomial, if used elsewhere, has its separate coefficientwise definition. The displayed unit ENDPOINTSOURCE12B6490A442BF168E00ED3A2SLOT83END proves monicity and retains the full complex phase. The transformed root radius is exactly bounded by

ENDPOINTSOURCE12B6490A442BF168E00ED3A2SLOT84END

For every ENDPOINTSOURCE12B6490A442BF168E00ED3A2SLOT85END, define the invertible linear map

ENDPOINTSOURCE12B6490A442BF168E00ED3A2SLOT86END ENDPOINTSOURCE12B6490A442BF168E00ED3A2SLOT87END

Substitution also defines an algebra isomorphism on the whole polynomial rings. In particular it sends the principal ideal ENDPOINTSOURCE12B6490A442BF168E00ED3A2SLOT88END onto ENDPOINTSOURCE12B6490A442BF168E00ED3A2SLOT89END, because ENDPOINTSOURCE12B6490A442BF168E00ED3A2SLOT90END and the factor ENDPOINTSOURCE12B6490A442BF168E00ED3A2SLOT91END is a nonzero scalar unit. Thus the induced quotient map is exactly

ENDPOINTSOURCE12B6490A442BF168E00ED3A2SLOT92END

with inverse induced by the second formula in (RB.26).

Apply (RB.26) to the original division ENDPOINTSOURCE12B6490A442BF168E00ED3A2SLOT93END. One obtains

ENDPOINTSOURCE12B6490A442BF168E00ED3A2SLOT94END

The last term has degree below ENDPOINTSOURCE12B6490A442BF168E00ED3A2SLOT95END. Uniqueness of division gives the exact commuting identity

ENDPOINTSOURCE12B6490A442BF168E00ED3A2SLOT96END

For raw jets, set ENDPOINTSOURCE12B6490A442BF168E00ED3A2SLOT97END. Let ENDPOINTSOURCE12B6490A442BF168E00ED3A2SLOT98END be defined using derivatives in ENDPOINTSOURCE12B6490A442BF168E00ED3A2SLOT99END at ENDPOINTSOURCE12B6490A442BF168E00ED3A2SLOT100END, and let ENDPOINTSOURCE12B6490A442BF168E00ED3A2SLOT101END use derivatives in ENDPOINTSOURCE12B6490A442BF168E00ED3A2SLOT102END at ENDPOINTSOURCE12B6490A442BF168E00ED3A2SLOT103END, with the same ordering and multiplicities. Repeated application of the ordinary chain rule gives, for every integer ENDPOINTSOURCE12B6490A442BF168E00ED3A2SLOT104END,

ENDPOINTSOURCE12B6490A442BF168E00ED3A2SLOT105END

Define the invertible diagonal map

ENDPOINTSOURCE12B6490A442BF168E00ED3A2SLOT106END

The complete raw-jet transport identity is therefore

ENDPOINTSOURCE12B6490A442BF168E00ED3A2SLOT107END

Every scalar ENDPOINTSOURCE12B6490A442BF168E00ED3A2SLOT108END, including its phase, appears in (RB.30)--(RB.31); the jet coordinates throughout are raw derivatives. Combining (RB.22), (RB.29), and (RB.31) proves preservation of the original raw jets as well as the transformed raw jets.

\subsubsection{6. Exact norm transfer back to the original coordinates}\label{endpoint-relation_balance_original--exact-norm-transfer-back-to-the-original-coordinates}

Define a norm on the original polynomial space by the displayed affine transport,

ENDPOINTSOURCE12B6490A442BF168E00ED3A2SLOT109END

This definition is independent of the choice of ENDPOINTSOURCE12B6490A442BF168E00ED3A2SLOT110END, because adding zero coefficients does not change the sum. Invertibility of ENDPOINTSOURCE12B6490A442BF168E00ED3A2SLOT111END proves that it is a norm. Equations (RB.19) and (RB.29), under the root-radius hypothesis in (RB.25), yield

ENDPOINTSOURCE12B6490A442BF168E00ED3A2SLOT112END

The same estimate expressed using the untransformed coefficient norm also has fully explicit constants. For ENDPOINTSOURCE12B6490A442BF168E00ED3A2SLOT113END, the binomial theorem gives

ENDPOINTSOURCE12B6490A442BF168E00ED3A2SLOT114END

The maximum-column characterization of the coefficient ENDPOINTSOURCE12B6490A442BF168E00ED3A2SLOT115END operator norm, proved after (RB.19), consequently gives the exact operator norms

ENDPOINTSOURCE12B6490A442BF168E00ED3A2SLOT116END

Thus the original-coordinate estimate is

ENDPOINTSOURCE12B6490A442BF168E00ED3A2SLOT117END

Finally, for any ENDPOINTSOURCE12B6490A442BF168E00ED3A2SLOT118END and any polynomial ENDPOINTSOURCE12B6490A442BF168E00ED3A2SLOT119END of degree below ENDPOINTSOURCE12B6490A442BF168E00ED3A2SLOT120END, the coefficient triangle inequality gives ENDPOINTSOURCE12B6490A442BF168E00ED3A2SLOT121END. Applying this to the exact representative in (RB.29) yields

ENDPOINTSOURCE12B6490A442BF168E00ED3A2SLOT122END

In particular the maximum on the segment ENDPOINTSOURCE12B6490A442BF168E00ED3A2SLOT123END, ENDPOINTSOURCE12B6490A442BF168E00ED3A2SLOT124END, is at most ENDPOINTSOURCE12B6490A442BF168E00ED3A2SLOT125END. Equations (RB.24), (RB.27), (RB.29), and (RB.31) specify the exact relation to the original annihilator, quotient, representative, and raw-jet coordinates; no measure or quotient metric is replaced in this coefficient calculation.
